% Part: proof-theory
% Chapter: proof-search
% Section: introduction

\documentclass[../../../include/open-logic-section]{subfiles}

\begin{document}

\olfileid[pt]{pt}{ps}{int}

\olsection{Introdução}

Os axiomas e as regras de inferência de sistemas de !!{proof}, como
\Log{G3c} ou \Log{N1i}, definem quais árvores de sequentes ou
!!{formula}s contam como !!{proof}s corretas. Dada uma árvore de
sequentes ou !!{formula}s desse tipo (possivelmente com informações
adicionais, como quais regras são aplicadas em cada passo, ou rótulos
para as premissas e para as regras de inferência que as descarregam),
as definições dos axiomas e das regras nos permitem verificar se a
árvore dada é uma !!{proof} correta. Um problema diferente, e em geral
mais difícil, é \emph{encontrar} !!a{proof} de um determinado sequente
ou de uma determinada !!{formula} a partir de premissas dadas. Esse é o
problema da \emph{busca de provas}.

Há um procedimento muito rudimentar para a busca de provas. Podemos
considerar as !!{formula}s sequências de símbolos de um alfabeto
finito. Sequentes e árvores de sequentes ou !!{formula}s também podem
ser considerados sequências de !!{formula}s (talvez usando colchetes
especiais para indicar as ramificações da árvore). Os axiomas e as
regras nos permitem verificar mecanicamente se uma sequência de
símbolos representa uma !!{proof} correta. Consequentemente, podemos
gerar mecanicamente \emph{todas as !!{proof}s corretas possíveis}
(gerando todas as sequências de símbolos do alfabeto usado para
representar !!{proof}s e verificando, para cada sequência candidata, se
ela representa uma !!{proof} correta). O procedimento rudimentar de
busca de provas é o seguinte: gerar todas as !!{proof}s corretas até
encontrarmos uma que seja uma !!{proof} do sequente dado (ou da
!!{formula} dada cujas premissas abertas estejam entre as premissas
dadas).

Um método melhor é usar o método ingênuo que você empregaria para
encontrar !!a{proof} à mão: começa-se pelo sequente final, escolhe-se
!!a{formula} e aplica-se uma regra de inferência ``de trás para a
frente'' para gerar uma premissa ou premissas que, se tivessem
!!{proof}s, combinariam com a última inferência para formar !!a{proof}
do sequente dado. Um método semelhante (embora mais complicado) é o que
você usaria para encontrar !!a{proof} em um sistema de dedução natural.

Mas esse método não é infalível: se você escolher a regra errada ou as
!!{formula}s na ordem errada, poderá chegar a um beco sem saída. Ele
também não está completamente especificado: em cada ponto, pode haver
mais de uma !!{formula} possível a escolher ou mais de uma regra que
possa ser aplicada de trás para a frente. Um \emph{algoritmo de busca
de provas} é um procedimento completamente especificado que encontrará
!!a{proof} caso exista uma (isto é, é infalível).

Alguns sistemas de !!{proof} permitem algoritmos de busca de provas
mais simples do que outros. Em sistemas de sequentes, o !!{main operator}
de !!a{formula}, e o fato de ela ocorrer à esquerda ou à
direita, determinam qual regra deve ser aplicada de trás para a frente.
Por exemplo, se você quiser encontrar !!a{proof} de $\Gamma \Sequent
\Delta, !A \land !B$ na qual $!A \land !B$ ocorre à direita, deverá
aplicar \RightR{\land} de trás para a frente e gerar as duas premissas
correspondentes $\Gamma \Sequent \Delta, !A$ e $\Gamma \Sequent
\Delta, !B$. Contudo, se seu sistema for \Log{G1c}, a presença da
contração torna as coisas mais difíceis porque oferece mais opções:
você também poderia aplicar \RightR{\Contraction} de trás para a frente
e gerar, em vez disso, a premissa $\Gamma \Sequent \Delta, !A \land
!B, !A \land !B$. E depois a premissa $\Gamma \Sequent \Delta, !A
\land !B, !A \land !B, !A \land !B$, e assim por diante. Em \Log{G2c},
a situação é ainda pior. Aplicar a regra \RightR{\land} de trás para a
frente também significa que você precisa adivinhar $\Gamma_1$,
$\Gamma_2$ tais que $\Gamma = \Gamma_1 \cup \Gamma_2$ e $\Delta_1$,
$\Delta_2$ tais que $\Delta = \Delta_1 \cup \Delta_2$, e tentar a
sorte com as premissas $\Gamma_1 \Sequent \Delta_1, !A$ e $\Gamma_2
\Sequent \Delta_2, !B$. Uma suposição errada pode conduzir a um beco
sem saída mais adiante, e você teria de voltar ao início e escolher
outros $\Gamma_1$, $\Gamma_2$, $\Delta_1$, $\Delta_2$. Se a !!{formula}
for $!A \lor !B$, você precisará escolher uma das duas possibilidades
para \RightR{\lor}, produzindo uma das duas premissas possíveis
$\Gamma \Sequent \Delta, !A$ ou $\Gamma \Sequent \Delta, !B$. Uma
escolha errada exigiria novamente voltar atrás.

O sistema \Log{G3c} não tem esses problemas. Ele não tem regras
estruturais, e a escolha da regra de inferência é determinada pelo
!!{main operator} e pelo lado em que a !!{formula} ocorre. Portanto,
\Log{G3c} é mais apropriado para a busca de provas do que \Log{G1c} ou
\Log{G2c}. Uma busca bem-sucedida produzirá !!a{proof} em \Log{G3c}, e
essa !!{proof} poderá então ser traduzida para \Log{G1c} ou \Log{G2c}
(ou para \Log{N1c}, aliás). Assim, ter um algoritmo de busca de provas
para \Log{G3c} e algoritmos de tradução de !!{proof}s de \Log{G3c} para
!!{proof}s em outros sistemas também resolve o problema da busca de
provas para esses sistemas.

Como sabemos que um algoritmo de busca de provas é infalível? Seria
necessário mostrar que, se o algoritmo não consegue encontrar
!!a{proof}, nenhuma !!{proof} pode existir. Afinal, se escolhermos um
algoritmo ruim, poderemos encontrar casos em que ele não encontra uma
prova mesmo quando existe uma. Esse problema não afeta o algoritmo
rudimentar, pois ele examina todas as !!{proof}s possíveis. Mas um
algoritmo de busca de provas deve ser eficiente e empregar o mínimo de
esforço.

A ideia essencial para mostrar que um algoritmo de busca de provas é
infalível é mostrar que, se o algoritmo falhar, poderemos mostrar, pela
maneira como ele falha, que o sequente não pode ter !!a{proof}. E a
maneira de mostrar isso é mostrar que ele não é válido. Na lógica
clássica de primeira ordem, os sequentes válidos $\Gamma \Sequent
\Delta$ são aqueles em que toda !!{structure} torna falsa pelo menos
uma !!{formula} de $\Gamma$ ou verdadeira pelo menos uma !!{formula}
de~$\Delta$. \Log{G3c} é \emph{correto}, isto é, !!{prove}s apenas
sequentes válidos. Isso é estabelecido por indução sobre !!{proof}s: os
axiomas são obviamente válidos e, se as premissas de uma regra forem
válidas, sua conclusão também será. Para mostrar que um algoritmo de
busca de provas é infalível, mostramos que, se a busca de !!a{proof} de
$\Gamma \Sequent \Delta$ falhar, poderemos definir !!a{structure} na
qual todas as !!{formula}s de $\Gamma$ sejam verdadeiras e todas as
!!{formula}s de~$\Delta$ sejam falsas. Isso também estabelece que
\Log{G3c} é \emph{completo}, isto é, que, se $\Gamma \Sequent \Delta$
for válido, haverá uma \Log{G3c}-!!{proof} dele. Como toda busca de
prova malsucedida mostra que $\Gamma \Sequent \Delta$ é inválido, toda
busca de prova de um sequente válido deve ter êxito.

\end{document}
